\clearpage
\begin{explanationbox}
\textbf{第一眼：两条抛物线为什么有两种加法？}
\begin{center}% p=1, a=2,b=-1. H=(x_B,y_A). Both signed displacements are negative.
\begin{tikzpicture}[font=\small,>=stealth]
\draw[gray!45,rounded corners=3pt] (0,0) rectangle (14.2,9.0);
\node at (7.1,8.45) {平方差怎样进入弦的方向：先横移，再纵移};
\begin{scope}[shift={(2.3,3.5)},scale=1.2]
\draw[->,gray] (-1.25,0)--(3.1,0) node[above] {$x$};
\draw[->,gray] (0,-1.55)--(0,3.0) node[left] {$y$};
\draw[teal!65!black,thick,samples=100,domain=-1.4:2.4,variable=\t] plot ({\t*\t/2},\t);
\draw[blue!75!black,thin] (0,0)--(2,2);
\draw[orange!85!black,thin] (0,0)--(.5,-1);
\draw[green!45!black,very thick] (2,2)--(.5,-1);
\draw[->,violet,thick] (2,2)--(.5,2);
\draw[->,orange!85!black,thick,dashed] (.5,2)--(.5,-1);
\draw[gray] (.5,1.8)--(.7,1.8)--(.7,2);
\node[violet] at (1.25,2.4) {$\Delta x$};
\node[orange!85!black] at (-.8,.6) {$\Delta y$};
\fill (0,0) circle(2pt) node[below left=7pt] {$O$};
\fill (2,2) circle(2pt) node[below right=8pt] {$A$};
\fill (.5,-1) circle(2pt) node[right=8pt] {$B$};
\fill (.5,2) circle(1.5pt) node[above=12pt] {$H$};
\end{scope}
\node[anchor=west,align=left] at (8.0,5.0)
 {图取 $p=1,\ a=2,\ b=-1$\\[12pt]
 $A(2,2),\quad B(\frac12,-1)$\\[15pt]
 $\Delta x=\dfrac{b^2-a^2}{2p}=-\dfrac32$\\[15pt]
 $\Delta y=b-a=-3$\\[15pt]
 $\dfrac{\Delta x}{\Delta y}=\dfrac{a+b}{2p}=\dfrac12$};
\node at (7.1,.6) {两个箭头分别指向左、下；横向抛物线所需的是有向比值 $\Delta x/\Delta y$。};
\end{tikzpicture}
\end{center}
图中的两次有向位移给出 $\Delta x/\Delta y$；平方差正好出现在横坐标差中。

抛物线 $y^2=2px$ 的\textbf{横坐标由纵坐标的平方决定}。若
\[
A=\left(\frac{a^2}{2p},a\right),\qquad
B=\left(\frac{b^2}{2p},b\right),
\]
则 $1/k_{OA}=x_A/y_A=a/(2p)$，$1/k_{OB}=b/(2p)$；而 $1/k_{AB}$ 恰好也等于 $(a+b)/(2p)$（只要 $AB$ 非竖直）。

换成 $x^2=2py$，\textbf{纵坐标由横坐标的平方决定}：
\[
A=\left(a,\frac{a^2}{2p}\right),\qquad
B=\left(b,\frac{b^2}{2p}\right),
\]
则 $k_{OA}=a/(2p)$，$k_{OB}=b/(2p)$，同时 $k_{AB}=(a+b)/(2p)$。

\textbf{核心原因只有一个：} $a^2-b^2=(a-b)(a+b)$，不是一个需要背诵的孤立巧合。

\end{explanationbox}
\clearpage
\begin{explanationbox}
\textbf{第二眼：把公式看成三角形三条边的方向关系。}
\begin{center}% p=1. Right panel is the exact x/y interchange of left panel.
\begin{tikzpicture}[font=\small,>=stealth,line join=round]
\draw[gray!45,rounded corners=3pt] (0,0) rectangle (14.2,10.5);
\node at (7.1,9.95) {看同一个三角形的三条边：蓝 $OA$，橙 $OB$，绿 $AB$};
\node at (3.55,9.05) {$y^2=2x$：横向开口};
\node at (10.65,9.05) {$x^2=2y$：纵向开口};
\begin{scope}[shift={(2.0,5.0)},scale=1.05]
\draw[->,gray] (-1.35,0)--(3.0,0) node[above] {$x$};
\draw[->,gray] (0,-1.65)--(0,3.0) node[left] {$y$};
\draw[teal!65!black,thick,samples=100,domain=-1.45:2.3,variable=\t] plot ({\t*\t/2},\t);
\fill[blue!4] (0,0)--(2,2)--(.5,-1)--cycle;
\draw[blue!75!black,very thick] (0,0)--(2,2);
\draw[orange!85!black,very thick] (0,0)--(.5,-1);
\draw[green!45!black,very thick] (2,2)--(.5,-1);
\fill (0,0) circle(2pt) node[below left=7pt] {$O$};
\fill (2,2) circle(2pt) node[below right=8pt] {$A$};
\fill (.5,-1) circle(2pt) node[right=8pt] {$B$};
\end{scope}
\begin{scope}[shift={(9.1,5.0)},scale=1.05]
\draw[->,gray] (-1.65,0)--(3.0,0) node[above] {$x$};
\draw[->,gray] (0,-1.35)--(0,3.0) node[left] {$y$};
\draw[teal!65!black,thick,samples=100,domain=-1.45:2.3,variable=\t] plot (\t,{\t*\t/2});
\fill[blue!4] (0,0)--(2,2)--(-1,.5)--cycle;
\draw[blue!75!black,very thick] (0,0)--(2,2);
\draw[orange!85!black,very thick] (0,0)--(-1,.5);
\draw[green!45!black,very thick] (2,2)--(-1,.5);
\fill (0,0) circle(2pt) node[below left=7pt] {$O$};
\fill (2,2) circle(2pt) node[below right=8pt] {$A$};
\fill (-1,.5) circle(2pt) node[above=14pt] {$B$};
\end{scope}
\node at (3.55,2.65) {$\textcolor{blue!75!black}{k_{OA}=1},\quad\textcolor{orange!85!black}{k_{OB}=-2},\quad\textcolor{green!45!black}{k_{AB}=2}$};
\node at (10.65,2.65) {$\textcolor{blue!75!black}{k_{OA}=1},\quad\textcolor{orange!85!black}{k_{OB}=-\frac12},\quad\textcolor{green!45!black}{k_{AB}=\frac12}$};
\node at (3.55,1.55) {$\displaystyle\frac11+\frac1{-2}=\frac12$};
\node at (10.65,1.55) {$\displaystyle 1+\left(-\frac12\right)=\frac12$};
\node at (7.1,.55) {两图同取 $p=1$，坐标互换且单位长度相同；三条边的颜色一一对应。};
\end{tikzpicture}
\end{center}

两条从顶点 $O$ 发出的射线 $OA,OB$ 一旦确定，弦 $AB$ 的方向不是随意的：
\[
\text{纵向开口： }k_{AB}=k_{OA}+k_{OB};
\qquad
\text{横向开口： }\frac1{k_{AB}}=\frac1{k_{OA}}+\frac1{k_{OB}}.
\]
蓝、橙两边确定后，绿色弦的方向随之确定。横向图读取 $\Delta x/\Delta y$；纵向图读取 $\Delta y/\Delta x$。

\end{explanationbox}
\clearpage
\begin{explanationbox}
\textbf{第三眼：对称弦是最值得记住的边界。}
\begin{center}
% p=1, free coordinates +1,-1, reflected with identical scale.
\begin{tikzpicture}[font=\small,>=stealth]
\draw[gray!45,rounded corners=3pt] (0,0) rectangle (14.2,10.0);
\node at (7.1,9.45) {对称的两个端点：竖直弦与水平弦的斜率含义不同};
\node at (3.55,8.45) {$y^2=2x$};
\node at (10.65,8.45) {$x^2=2y$};
\begin{scope}[shift={(2.6,5.1)},scale=1.25]
\draw[->,gray] (-1.5,0)--(2.1,0) node[above] {$x$};
\draw[->,gray] (0,-1.5)--(0,2.1) node[left] {$y$};
\draw[teal!65!black,thick,samples=90,domain=-1.4:1.8,variable=\t] plot ({\t*\t/2},\t);
\draw[blue!75!black,thick] (0,0)--(.5,1);
\draw[orange!85!black,thick] (0,0)--(.5,-1);
\draw[green!45!black,very thick] (.5,1)--(.5,-1);
\fill (0,0) circle(2pt) node[below left=7pt] {$O$};
\fill (.5,1) circle(2pt) node[right=9pt] {$A$};
\fill (.5,-1) circle(2pt) node[right=9pt] {$B$};
\end{scope}
\begin{scope}[shift={(9.7,5.1)},scale=1.25]
\draw[->,gray] (-1.5,0)--(2.1,0) node[above] {$x$};
\draw[->,gray] (0,-1.5)--(0,2.1) node[left] {$y$};
\draw[teal!65!black,thick,samples=90,domain=-1.4:1.8,variable=\t] plot (\t,{\t*\t/2});
\draw[blue!75!black,thick] (0,0)--(1,.5);
\draw[orange!85!black,thick] (0,0)--(-1,.5);
\draw[green!45!black,very thick] (-1,.5)--(1,.5);
\fill (0,0) circle(2pt) node[below left=7pt] {$O$};
\fill (1,.5) circle(2pt) node[above=14pt] {$A$};
\fill (-1,.5) circle(2pt) node[above=14pt] {$B$};
\end{scope}
\node at (3.55,2.6) {$AB:\ x=\frac12$，$k_{AB}$ 不存在};
\node at (10.65,2.6) {$AB:\ y=\frac12$，$k_{AB}=0$};
\node at (3.55,1.55) {$\displaystyle\frac{x_B-x_A}{y_B-y_A}=0$};
\node at (10.65,1.55) {$\displaystyle k_{OA}+k_{OB}=\frac12-\frac12=0$};
\node at (7.1,.55) {同一尺度下，横向模型的竖直弦关于 $y=x$ 翻转后成为纵向模型的水平弦。};
\end{tikzpicture}

\end{center}
两点关于轴对称时，横向开口的弦\textbf{竖直}，普通 $k_{AB}$ 不存在；纵向开口的弦\textbf{水平}，$k_{AB}=0$。这两种情况不能混淆。
\end{explanationbox}

\clearpage
\begin{explanationbox}
\textbf{直接推广一：为什么顶点不在原点也成立？}
\begin{center}% Exact translation: Y=X^2 -> y=(x-2)^2+1; equal x/y units.
\begin{tikzpicture}[font=\small,>=stealth]
\draw[gray!45,rounded corners=3pt] (0,0) rectangle (14.2,10.0);
\node at (7.1,9.45) {只平移图形，三条边的方向都不变};
\node at (3.55,8.55) {$Y=X^2$};
\node at (10.65,8.55) {$y=(x-2)^2+1$};
\begin{scope}[shift={(3.2,3.8)},scale=1.15]
\draw[->,gray] (-2.1,0)--(2.1,0) node[above] {$X$};
\draw[->,gray] (0,-.75)--(0,3.4) node[left] {$Y$};
\draw[teal!65!black,thick,samples=80,domain=-1.7:1.7,variable=\t] plot (\t,{\t*\t});
\draw[blue!75!black,very thick] (0,0)--(1,1);
\draw[orange!85!black,very thick] (0,0)--(-1.5,2.25);
\draw[green!45!black,very thick] (1,1)--(-1.5,2.25);
\fill (0,0) circle(2pt) node[below left=8pt] {$V$};
\fill (1,1) circle(2pt) node[right=12pt] {$A$};
\fill (-1.5,2.25) circle(2pt) node[right=14pt] {$B$};
\end{scope}
\begin{scope}[shift={(8.1,2.65)},scale=1.15]
\draw[->,gray] (-.25,0)--(4.5,0) node[above] {$x$};
\draw[->,gray] (0,-.3)--(0,4.4) node[left] {$y$};
\draw[teal!65!black,thick,samples=80,domain=.3:3.7,variable=\t] plot (\t,{(\t-2)*(\t-2)+1});
\draw[blue!75!black,very thick] (2,1)--(3,2);
\draw[orange!85!black,very thick] (2,1)--(.5,3.25);
\draw[green!45!black,very thick] (3,2)--(.5,3.25);
\draw[gray!55,dashed] (2,0)--(2,1)--(0,1);
\fill (2,1) circle(2pt) node[below right=10pt] {$V(2,1)$};
\fill (3,2) circle(2pt) node[right=12pt] {$A$};
\fill (.5,3.25) circle(2pt) node[right=14pt] {$B$};
\end{scope}
\node at (3.55,1.8) {$k_{VA}=1,\quad k_{VB}=-\frac32$};
\node at (10.65,1.8) {$k_{VA}=1,\quad k_{VB}=-\frac32$};
\node at (7.1,.65) {两图同尺度：横、纵坐标差保持不变，故 $k_{AB}=1-\frac32=-\frac12$。};
\end{tikzpicture}
\end{center}
把左图向右移 $2$、向上移 $1$，顶点和两个端点一起移动。
每条边两端的横坐标都加 $2$、纵坐标都加 $1$，所以相减时新增常数抵消，斜率不变。

一般地，顶点为 $V(h,k)$ 的二次函数可写为 $y=a(x-h)^2+k$。
令 $X=x-h,\ Y=y-k$，得到 $Y=aX^2$。顶点连线斜率分别是 $aX_1,aX_2$，弦斜率是 $a(X_1+X_2)$。
\textbf{规律来自平方差，顶点的位置不影响它。}
\end{explanationbox}

\clearpage
\begin{explanationbox}
\textbf{直接推广二：同一端点，两种直线为什么差一个 $2$？}
\begin{center}% x^2=2y, p=1; A=(-1,1/2), B=(2,2); slopes -1,2,1/2.
\begin{tikzpicture}[font=\small,>=stealth]
\draw[gray!45,rounded corners=3pt] (0,0) rectangle (14.2,10.0);
\node at (7.1,9.45) {纵向：两端切线的斜率取平均，得到弦的斜率};
\begin{scope}[shift={(3.8,3.1)},scale=1.15]
\draw[->,gray] (-2.4,0)--(3.4,0) node[above] {$x$};
\draw[->,gray] (0,-1.3)--(0,4.3) node[left] {$y$};
\draw[teal!65!black,thick,samples=100,domain=-2.2:2.8,variable=\t] plot (\t,{\t*\t/2});
\draw[blue!75!black,dashed,thick] (0,0)--(-1,.5);
\draw[orange!85!black,dashed,thick] (0,0)--(2,2);
\draw[blue!75!black,very thick] (-2.2,1.7)--(.6,-1.1);
\draw[orange!85!black,very thick] (.6,-.8)--(2.8,3.6);
\draw[green!45!black,very thick] (-1,.5)--(2,2);
\fill (0,0) circle(2pt) node[below left=10pt] {$O$};
\fill (-1,.5) circle(2pt) node[above=14pt] {$A$};
\fill (2,2) circle(2pt) node[right=14pt] {$B$};
\node[blue!75!black] at (-1.6,2.1) {$t_A$};
\node[orange!85!black] at (3.25,3.25) {$t_B$};
\end{scope}
\node[anchor=west,align=left] at (8.9,5.7)
 {$x^2=2y,\quad p=1$\\[13pt]
 $A(-1,\frac12),\ B(2,2)$\\[15pt]
 $k_A=-1,\quad k_B=2$\\[15pt]
 $k_{AB}=\dfrac{-1+2}{2}=\dfrac12$};
\node at (7.1,1.3) {虚线是顶点连线，实线 $t_A,t_B$ 是切线；二者不能混淆。};
\node at (7.1,.6) {$k_{OA}=-\frac12,\ k_{OB}=1$，切线斜率分别是顶点连线斜率的 $2$ 倍。};
\end{tikzpicture}
\end{center}
图中虚线连接顶点与端点，实线在端点处与曲线相切。对 $x^2=2py$ 上的非顶点 $P(u,u^2/(2p))$，$u\ne0$，顶点连线斜率为 $u/(2p)$，切线斜率为 $u/p$。

因此两条切线斜率之和，是两条顶点连线斜率之和的 $2$ 倍。\textbf{弦斜率要取切线斜率的平均值}，不能把切线斜率直接相加。
已知两端切线方向时，无须分别求出两端点，也能算出弦的方向。
\end{explanationbox}

\clearpage
\begin{explanationbox}
\textbf{横向为什么改成倒斜率取平均？}
\begin{center}% Coordinate interchange of the vertical tangent diagram; identical units.
\begin{tikzpicture}[font=\small,>=stealth]
\draw[gray!45,rounded corners=3pt] (0,0) rectangle (14.2,10.0);
\node at (7.1,9.45) {横向：坐标互换后，改为两端切线的倒斜率取平均};
\begin{scope}[shift={(2.6,4.4)},scale=1.15]
\draw[->,gray] (-1.3,0)--(4.3,0) node[above] {$x$};
\draw[->,gray] (0,-2.4)--(0,3.4) node[left] {$y$};
\draw[teal!65!black,thick,samples=100,domain=-2.2:2.8,variable=\t] plot ({\t*\t/2},\t);
\draw[blue!75!black,dashed,thick] (0,0)--(.5,-1);
\draw[orange!85!black,dashed,thick] (0,0)--(2,2);
\draw[blue!75!black,very thick] (1.7,-2.2)--(-1.1,.6);
\draw[orange!85!black,very thick] (-.8,.6)--(3.6,2.8);
\draw[green!45!black,very thick] (.5,-1)--(2,2);
\fill (0,0) circle(2pt) node[below left=12pt] {$O$};
\fill (.5,-1) circle(2pt) node[right=14pt] {$A$};
\fill (2,2) circle(2pt) node[above=15pt] {$B$};
\node[blue!75!black] at (2.8,-1.7) {$t_A$};
\node[orange!85!black] at (3.6,1.95) {$t_B$};
\end{scope}
\node[anchor=west,align=left] at (8.9,5.7)
 {$y^2=2x,\quad p=1$\\[13pt]
 $A(\frac12,-1),\ B(2,2)$\\[15pt]
 $k_A=-1,\quad k_B=\frac12$\\[15pt]
 $\dfrac1{k_{AB}}=\dfrac{-1+2}{2}=\dfrac12$};
\node at (7.1,1.3) {同一对照模型交换 $x,y$；绿色弦的普通斜率变为 $k_{AB}=2$。};
\node at (7.1,.6) {这里的 $A,B$ 都不是顶点，弦不竖直，所有倒斜率才可按实数计算。};
\end{tikzpicture}
\end{center}
与纵向图交换横、纵坐标后，非零有限斜率变成倒数。也可直接读横向位移：对非顶点 $P(u^2/(2p),u)$，$u\ne0$，顶点连线的倒斜率是 $u/(2p)$，切线的倒斜率是 $u/p$。

图中 $\frac12(-1+2)=\frac12$ 给出的是\textbf{弦的倒斜率}，所以普通弦斜率为 $2$。
若两端纵坐标互为相反数，平均值为 $0$，实际弦竖直；此时只能用坐标差说明，不能写成普通斜率的倒数。
\end{explanationbox}

\clearpage
\begin{explanationbox}
\textbf{核心技巧：看到倒数相加，就想到“乘积除以和”。}
\begin{center}% p=1, free coordinates +1,-1, reflected with identical scale.
\begin{tikzpicture}[font=\small,>=stealth]
\draw[gray!45,rounded corners=3pt] (0,0) rectangle (14.2,10.0);
\node at (7.1,9.45) {对称的两个端点：竖直弦与水平弦的斜率含义不同};
\node at (3.55,8.45) {$y^2=2x$};
\node at (10.65,8.45) {$x^2=2y$};
\begin{scope}[shift={(2.6,5.1)},scale=1.25]
\draw[->,gray] (-1.5,0)--(2.1,0) node[above] {$x$};
\draw[->,gray] (0,-1.5)--(0,2.1) node[left] {$y$};
\draw[teal!65!black,thick,samples=90,domain=-1.4:1.8,variable=\t] plot ({\t*\t/2},\t);
\draw[blue!75!black,thick] (0,0)--(.5,1);
\draw[orange!85!black,thick] (0,0)--(.5,-1);
\draw[green!45!black,very thick] (.5,1)--(.5,-1);
\fill (0,0) circle(2pt) node[below left=7pt] {$O$};
\fill (.5,1) circle(2pt) node[right=9pt] {$A$};
\fill (.5,-1) circle(2pt) node[right=9pt] {$B$};
\end{scope}
\begin{scope}[shift={(9.7,5.1)},scale=1.25]
\draw[->,gray] (-1.5,0)--(2.1,0) node[above] {$x$};
\draw[->,gray] (0,-1.5)--(0,2.1) node[left] {$y$};
\draw[teal!65!black,thick,samples=90,domain=-1.4:1.8,variable=\t] plot (\t,{\t*\t/2});
\draw[blue!75!black,thick] (0,0)--(1,.5);
\draw[orange!85!black,thick] (0,0)--(-1,.5);
\draw[green!45!black,very thick] (-1,.5)--(1,.5);
\fill (0,0) circle(2pt) node[below left=7pt] {$O$};
\fill (1,.5) circle(2pt) node[above=14pt] {$A$};
\fill (-1,.5) circle(2pt) node[above=14pt] {$B$};
\end{scope}
\node at (3.55,2.6) {$AB:\ x=\frac12$，$k_{AB}$ 不存在};
\node at (10.65,2.6) {$AB:\ y=\frac12$，$k_{AB}=0$};
\node at (3.55,1.55) {$\displaystyle\frac{x_B-x_A}{y_B-y_A}=0$};
\node at (10.65,1.55) {$\displaystyle k_{OA}+k_{OB}=\frac12-\frac12=0$};
\node at (7.1,.55) {同一尺度下，横向模型的竖直弦关于 $y=x$ 翻转后成为纵向模型的水平弦。};
\end{tikzpicture}
\end{center}
对非零实数 $u,v$，通分恒有
\[
\frac1u+\frac1v=\frac{u+v}{uv}.
\]
若它还等于 $1/w$（$w\ne0$），则 $u+v\ne0$，取倒数得
\[
\boxed{w=\frac{uv}{u+v}}.
\]
\textbf{正例：}$u=2,v=6$ 时，$w=12/8=3/2$。
\textbf{边界：}$u=2,v=-2$ 时，倒数和为 $0$，没有有限非零实数 $w$ 满足 $1/w=0$；用于横向抛物线，就是图中的竖直弦。

\textbf{辨析：}二次函数平移和切线平均关系来自同一个平方差结构；焦点弦倒数和来自焦点位置与距离公式，直角三角形倒数平方和来自面积与勾股。它们只有代数形式相似，不能互相代入。
\end{explanationbox}
